\textbf{Phần 1. Trắc nghiệm lựa chọn}
\setcounterpageref{ex}{0}

\begin{ex} %[3D1V04]
Tính giá trị của biểu thức: $A = 2 \times 8 + 5 \times 6 - 20$.
\choice
{$24$}
{$28$}
{$30$}
{\True $26$}
\loigiai{
Thực hiện các phép tính nhân trước, cộng trừ sau:\\
$$A = 16 + 30 - 20 = 46 - 20 = 26$$
\textbf{Đáp án của câu hỏi là D}
}
\end{ex}

\begin{ex} %[3D1V04]
Cho hai biểu thức $M = 5 \times 7 - 12$ và $N = 2 \times 9 + 6$. So sánh giá trị của $M$ và $N$:
\choice
{$M > N$}
{\True $M < N$}
{$M = N$}
{$M + N = 50$}
\loigiai{
Tính giá trị của từng biểu thức:\\
- $M = 5 \times 7 - 12 = 35 - 12 = 23$\\
- $N = 2 \times 9 + 6 = 18 + 6 = 24$\\
Vì $23 < 24$ nên $M < N$.\\
\textbf{Đáp án của câu hỏi là B}
}
\end{ex}

\begin{ex} %[3D1V04]
Tìm giá trị của $x$, biết: $5 \times x + 25 = 60$.
\choice
{$5$}
{$6$}
{\True $7$}
{$8$}
\loigiai{
Xem $5 \times x$ là số hạng chưa biết:\\
$$5 \times x = 60 - 25$$
$$5 \times x = 35$$
$$x = 35 : 5 = 7$$
\textbf{Đáp án của câu hỏi là C}
}
\end{ex}

\begin{ex} %[3D1V04]
Một nông trại nuôi $9$ con gà và $6$ con vịt. Hỏi đàn gà và đàn vịt đó có tất cả bao nhiêu cái chân?
\choice
{\True $30$ cái chân}
{$28$ cái chân}
{$32$ cái chân}
{$26$ cái chân}
\loigiai{
Mỗi con gà và mỗi con vịt đều có $2$ cái chân.\\
Tổng số con gà và vịt là:
$$9 + 6 = 15\text{ (con)}$$
Đàn gà và vịt có tất cả số cái chân là:
$$15 \times 2 = 30\text{ (cái chân)}$$
\textbf{Đáp án của câu hỏi là A}
}
\end{ex}

\begin{ex} %[3D1V04]
Cô giáo có $45$ cái kẹo, cô chia đều cho các bạn học sinh giỏi, mỗi bạn được $5$ cái kẹo. Sau đó cô giáo lại thưởng thêm cho mỗi bạn $2$ cái kẹo nữa. Hỏi cô giáo đã thưởng thêm tất cả bao nhiêu cái kẹo?
\choice
{$14$ cái kẹo}
{$16$ cái kẹo}
{$20$ cái kẹo}
{\True $18$ cái kẹo}
\loigiai{
Số bạn học sinh giỏi được nhận kẹo là:
$$45 : 5 = 9\text{ (bạn)}$$
Cô giáo đã thưởng thêm tất cả số cái kẹo là:
$$9 \times 2 = 18\text{ (cái kẹo)}$$
\textbf{Đáp án của câu hỏi là D}
}
\end{ex}

\begin{ex} %[3D1V04]
Cho dãy số: $5; 10; 15; 20; \dots; 50$. Tổng của số hạng thứ tư và số hạng thứ chín trong dãy số trên là
\choice
{$60$}
{\True $65$}
{$70$}
{$55$}
\loigiai{
Các số trong dãy là tích của $5$ với số thứ tự của nó:\\
- Số hạng thứ tư là: $5 \times 4 = 20$\\
- Số hạng thứ chín là: $5 \times 9 = 45$\\
Tổng hai số hạng là:
$$20 + 45 = 65$$
\textbf{Đáp án của câu hỏi là B}
}
\end{ex}

\begin{ex} %[3D1V04]
Một số tự nhiên khi chia cho $5$ thì được thương là $8$ và số dư là số dư lớn nhất có thể. Số tự nhiên đó là
\choice
{$42$}
{$43$}
{\True $44$}
{$45$}
\loigiai{
Vì số chia là $5$ nên số dư lớn nhất có thể có là $4$.\\
Số tự nhiên cần tìm là:
$$5 \times 8 + 4 = 40 + 4 = 44$$
\textbf{Đáp án của câu hỏi là C}
}
\end{ex}

\begin{ex} %[3D1V04]
Có 5 túi quà, mỗi túi gồm $4$ viên kẹo cam và $2$ viên kẹo dâu. Hỏi trong 5 túi quà đó có tất cả bao nhiêu viên kẹo?
\choice
{\True $30$ viên kẹo}
{$25$ viên kẹo}
{$35$ viên kẹo}
{$28$ viên kẹo}
\loigiai{
Mỗi túi có số viên kẹo là:
$$4 + 2 = 6\text{ (viên kẹo)}$$
Trong 5 túi có tất cả số viên kẹo là:
$$6 \times 5 = 30\text{ (viên kẹo)}$$
\textbf{Đáp án của câu hỏi là A}
}
\end{ex}

\begin{ex} %[3D1V04]
Tìm giá trị của $x$, biết: $x : 2 + 15 = 24$.
\choice
{$16$}
{$20$}
{$14$}
{\True $18$}
\loigiai{
Xem $x : 2$ là số hạng chưa biết:\\
$$x : 2 = 24 - 15$$
$$x : 2 = 9$$
$$x = 9 \times 2 = 18$$
\textbf{Đáp án của câu hỏi là D}
}
\end{ex}

\begin{ex} %[3D1V04]
Bình nghĩ ra một số. Bình lấy số đó nhân với $5$ rồi trừ đi $18$ thì được kết quả là $27$. Số mà Bình đã nghĩ ra là
\choice
{$8$}
{\True $9$}
{$7$}
{$10$}
\loigiai{
Gọi số Bình nghĩ là $x$. Ta có:\\
$$x \times 5 - 18 = 27$$
$$x \times 5 = 27 + 18$$
$$x \times 5 = 45$$
$$x = 45 : 5 = 9$$
\textbf{Đáp án của câu hỏi là B}
}
\end{ex}

\begin{ex} %[3D1V04]
Một cửa hàng có $50\text{ kg}$ đường. Người ta đóng đều vào $6$ túi lớn, mỗi túi $5\text{ kg}$ đường. Số đường còn lại đóng đều vào các túi nhỏ, mỗi túi chứa $2\text{ kg}$ đường. Hỏi người ta đóng được bao nhiêu túi nhỏ?
\choice
{$8$ túi}
{$9$ túi}
{\True $10$ túi}
{$12$ túi}
\loigiai{
Số đường đóng vào 6 túi lớn là:
$$6 \times 5 = 30\text{ (kg)}$$
Số ki-lô-gam đường còn lại là:
$$50 - 30 = 20\text{ (kg)}$$
Số túi nhỏ đóng được là:
$$20 : 2 = 10\text{ (túi)}$$
\textbf{Đáp án của câu hỏi là C}
}
\end{ex}

\begin{ex} %[3D1V04]
Tính nhanh giá trị của biểu thức: $5 \times 3 + 5 \times 4 + 5 \times 3$.
\choice
{\True $50$}
{$45$}
{$55$}
{$60$}
\loigiai{
Áp dụng tính chất phân phối của phép nhân đối với phép cộng:\\
$$5 \times 3 + 5 \times 4 + 5 \times 3 = 5 \times (3 + 4 + 3) = 5 \times 10 = 50$$
\textbf{Đáp án của câu hỏi là A}
}
\end{ex}

\textbf{Phần 2. Trắc nghiệm đúng sai}
\setcounterpageref{ex}{0}

\begin{ex} %[3D1V04]
Xét tính đúng sai của các nhận định sau về bảng nhân và bảng chia 2, 5:
\choiceTF
{\True Tích của một số tự nhiên bất kì nhân với $2$ luôn luôn là một số chẵn.}
{Một số chia hết cho $5$ thì luôn luôn có chữ số tận cùng là $5$.}
{\True Nếu $x : 5 = 7$ thì giá trị của $x$ bằng $35$.}
{\True Giá trị của biểu thức $2 \times 5$ bằng giá trị của biểu thức $5 \times 2$.}
\loigiai{
\begin{itemize}
    \item Ý a) Đúng: Bất kì số tự nhiên nào nhân với 2 cũng được một số chia hết cho 2 (số chẵn).
    \item Ý b) Sai: Một số chia hết cho 5 có chữ số tận cùng là 0 hoặc 5 (ví dụ 10, 20 chia hết cho 5 nhưng tận cùng là 0).
    \item Ý c) Đúng: $x = 7 \times 5 = 35$.
    \item Ý d) Đúng: Phép nhân có tính chất giao hoán ($2 \times 5 = 5 \times 2 = 10$).
\end{itemize}
\begin{itemize}
    \item a) \textbf{Đúng}.
    \item b) \textbf{Sai}.
    \item c) \textbf{Đúng}.
    \item d) \textbf{Đúng}.
\end{itemize}
}
\end{ex}

\begin{ex} %[3D1V04]
Một nhóm có $10$ bạn học sinh xếp thành $2$ hàng đều nhau, mỗi bạn cầm trên tay $5$ quả bóng bay. Xét tính đúng sai:
\choiceTF
{\True Mỗi hàng có $5$ bạn học sinh.}
{\True Cả nhóm học sinh có tất cả $50$ quả bóng bay.}
{Mỗi hàng học sinh cầm tất cả $20$ quả bóng bay.}
{\True Nếu mỗi bạn thả bay đi $1$ quả bóng thì cả nhóm còn lại $40$ quả bóng bay.}
\loigiai{
\begin{itemize}
    \item Ý a) Đúng: Số bạn mỗi hàng là: $10 : 2 = 5$ (bạn).
    \item Ý b) Đúng: Cả nhóm có: $10 \times 5 = 50$ (quả bóng bay).
    \item Ý c) Sai: Mỗi hàng có 5 bạn nên cầm số bóng là: $5 \times 5 = 25$ (quả bóng bay).
    \item Ý d) Đúng: Mỗi bạn còn 4 quả bóng, 10 bạn còn: $10 \times 4 = 40$ (quả bóng bay).
\end{itemize}
\begin{itemize}
    \item a) \textbf{Đúng}.
    \item b) \textbf{Đúng}.
    \item c) \textbf{Sai}.
    \item d) \textbf{Đúng}.
\end{itemize}
}
\end{ex}

\textbf{Phần 3. Trả lời ngắn}
\setcounterpageref{ex}{0}

\begin{ex} %[3D1V04]
Tính giá trị của biểu thức: $5 \times 9 - 2 \times 8$.
\shortans{29}
\loigiai{
\textbf{Đáp số:} $29$\\
Thực hiện phép nhân trước:\\
$$5 \times 9 - 2 \times 8 = 45 - 16 = 29$$
}
\end{ex}

\begin{ex} %[3D1V04]
Một người nuôi $8$ con thỏ. Hỏi $8$ con thỏ đó có tất cả bao nhiêu cái tai?
\shortans{16}
\loigiai{
\textbf{Đáp số:} $16$\\
Mỗi con thỏ có 2 cái tai. Số cái tai của 8 con thỏ là:
$$8 \times 2 = 16\text{ (cái tai)}$$
}
\end{ex}

\begin{ex} %[3D1V04]
Tìm số bị chia trong một phép chia cho $5$, biết thương bằng $6$ và số dư bằng $3$.
\shortans{33}
\loigiai{
\textbf{Đáp số:} $33$\\
Số bị chia là:
$$5 \times 6 + 3 = 30 + 3 = 33$$
}
\end{ex}

\begin{ex} %[3D1V04]
Lớp 3A có $35$ học sinh. Cô giáo chia đều học sinh thành các nhóm, mỗi nhóm có $5$ bạn. Sau đó mỗi nhóm cử $2$ bạn tham gia đội sao đỏ. Hỏi lớp 3A có tất cả bao nhiêu bạn tham gia đội sao đỏ?
\shortans{14}
\loigiai{
\textbf{Đáp số:} $14$\\
Lớp 3A chia được số nhóm là:
$$35 : 5 = 7\text{ (nhóm)}$$
Số bạn tham gia đội sao đỏ là:
$$7 \times 2 = 14\text{ (bạn)}$$
}
\end{ex}

\begin{ex} %[3D1V04]
Tìm giá trị của $x$, biết: $x : 5 = 14 - 6$.
\shortans{40}
\loigiai{
\textbf{Đáp số:} $40$\\
Tính vế phải:\\
$$x : 5 = 8$$
$$x = 8 \times 5 = 40$$
}
\end{ex}

\begin{ex} %[3D1V04]
Tính nhanh giá trị của biểu thức: $2 \times 1 + 2 \times 2 + 2 \times 3 + 2 \times 4$.
\shortans{20}
\loigiai{
\textbf{Đáp số:} $20$\\
Đặt thừa số chung:\\
$$2 \times (1 + 2 + 3 + 4) = 2 \times 10 = 20$$
}
\end{ex}
